\cvsection{Professional experience}
\vspace{8pt}

\cvsubsection{Positions}
%-------------------------------------------------------------------------------
%	Content
%-------------------------------------------------------------------------------
\begin{cventries}
	
	%---------------------------------------------------------
    \cventry
	{European Southern Observatory} % Degree - Changed it to institution so it is not highlighted
	{Postdoctoral Fellow} % Institution -  Changed it to degree so it is highlighted
	{September 2024 - Present} % Location - Changed to dates because I want this to be in color
	{Garching, Germany} % Date(s) - Changed to location because it is not as important 
	{ \begin{cvitems} % Description(s) bullet points
			\item{Study the effects of asymmetric infall in the development of protostellar and protoplanetary disks.}
            \item{Gather a sample of streamers toward protostars with different levels of accretion, with the goal of studying the relation between streamer mass delivery and protostellar growth.}
            \item{Functional duties at the European ALMA Regional Center (ARC) at ESO Headquarters}
		\end{cvitems}
	}
	\cventry
	{Max Planck Institute for Extraterrestrial Physics and Ludwig-Maximilians-Universit\"at} % Degree - Changed it to institution so it is not highlighted
	{Postdoctoral Researcher} % Institution -  Changed it to degree so it is highlighted
	{May 2024 - August 2024} % Location - Changed to dates because I want this to be in color
	{Garching, Germany} % Date(s) - Changed to location because it is not as important 
	{ \begin{cvitems} % Description(s) bullet points
			\item{Part of the data calibration and analysis for the NOEMA large program PRODIGE (PI: P. Caselli and Th. Henning)}
			\item{Analysis of the kinematics of protostellar envelopes using interferometers (e.g. ALMA and NOEMA) and single-dish telescopes (e.g. IRAM 30m)}
            \item{Short bridging contract between end of PhD and start of Fellowship}
		\end{cvitems}
	}
	
	
	%---------------------------------------------------------
\end{cventries}
\vspace{10pt}
\cvsubsection{Professional Activities}

\paragraphstyle
\vspace{8pt}

\begin{cvitems}
    \item \textbf{Telescopes}: PI of accepted projects in ALMA, NOEMA, APEX, IRAM 30m, Green Bank Telescope (GBT). Observing experience in APEX, IRAM 30m and GBT. 
    \item \textbf{Talks}: International conferences, e.g. European Astronomical Society Annual meeting (2021 and 2022) and the Latin American Regional IAU Meeting (2019). Contributed talks in conferences, e.g. Early Phases of Star Formation (Germany, 2024) and The promises and challenges of the ALMA Wideband Sensitivity Upgrade (ESO, 2024).
    \item \textbf{Seminars}: Invited to give seminars such as in Cardiff Astro Seminar Series (Cardiff, Wales, 2024), Seminarios de Fisíca Universidad de Santiago de Chile (Chile, 2023)
    \item \textbf{Outreach}: participated in several public outreach activities, in particular focused on improving visibility of women in STEM fields, such as the projects "Soapbox Science" (Munich, Germany) and "Cazadoras de Estrellas" (Santiago, Chile). 
    \item \textbf{Teaching}: experience as Teaching Assistance for astronomy and physics courses during undergrad and master studies (Universidad de Chile).
\end{cvitems}

% \begin{cvskills}
	
% 	%---------------------------------------------------------
% 	\cvskill
% 	{Telescopes} % Category
% 	{PI of accepted projects in ALMA, NOEMA, APEX, IRAM 30m, Green Bank Telescope (GBT).
%     \newline Observing experience in APEX, IRAM 30m and GBT. }% Skills
	
% 	\cvskill{Talks}
% {International conferences, e.g. European Astronomical Society Annual meeting (2021 and 2022) and the Latin American Regional IAU Meeting (2019).\newline Contributed talks in conferences, e.g. Early Phases of Star Formation (Germany, 2024) and \newline The promises and challenges of the ALMA Wideband Sensitivity Upgrade (ESO, 2024).}

% \cvskill{Seminars}
% {Invited to give seminars such as in Cardiff Astro Seminar Series (Cardiff, Wales, 2024), Seminarios de Fisíca Universidad de Santiago de Chile (Chile, 2023)}

% \end{cvskills}